\documentclass[11pt,letterpaper]{article}
\usepackage[T1]{fontenc}
\usepackage[utf8]{inputenc}
\usepackage[scaled=0.95]{helvet}
\renewcommand{\familydefault}{\sfdefault}
\usepackage[top=1.9cm,bottom=1.6cm,left=2.3cm,right=2.3cm]{geometry}
\usepackage{microtype}
\usepackage[hidelinks]{hyperref}
\pagestyle{empty}
\setlength{\parindent}{0pt}
\setlength{\parskip}{8pt}

\begin{document}

{\large\bfseries Grover Aruquipa}\par\vspace{-6pt}
{\small Qu\'ebec, QC, Canada \ \textbullet\ \ grover.aruquipa.9@gmail.com \ \textbullet\ \ groveraruquipa.github.io}

\vspace{4pt}
Haply Robotics, Montr\'eal\\
Physical AI Researcher

Dear Haply team,

I am applying for the Physical AI Researcher position. I am finishing my PhD in Mechanical Engineering at Universit\'e Laval with Prof.\ Cl\'ement Gosselin and Prof.\ Patrice Lambert, and I defend in early 2027. What I like about Haply is the idea that a robot should learn from touch and force, and not only from images. Force is also at the center of my own work.

During my PhD I designed, built and controlled three manipulators with integrated grippers, with 4, 7 and 8 actuators. I integrated the motors, drives and encoders, wrote the real-time control at 500\,Hz, and spent many hours on the bench making the real robot match the model. I write the URDF and MJCF files of my robots myself and validate them in MuJoCo before anything runs on hardware. My current robot is a Delta with an integrated Sarrus gripper, designed to sense the grasp force through the motor currents. I am training a reinforcement learning policy for it in NVIDIA Isaac Sim, and taking that policy from simulation to the real robot is my next step, which is exactly the kind of work your team describes.

I program in Python and C++, I have used PyTorch, and before my PhD I worked on computer vision for industrial anomaly detection at Sencon and on machine learning for sensor data at Hitachi ABB Power Grids R\&D in Zurich. I enjoy reading new papers and turning them into something that runs on a real robot, and I like working with people from hardware and software at the same time. I supervise two graduate students day to day, and I led the Laval team that won first place in the ASME Student Mechanisms and Robotics Design Competition in 2025.

I am available from January 2027 and ready to move to Montr\'eal. I will be eligible for an open Post-Graduation Work Permit, so I do not need sponsorship. You can see my robots on my website, groveraruquipa.github.io.

Thank you for your time,

Grover Aruquipa

\end{document}
